\documentclass[11pt]{article}
\usepackage{amsmath,amssymb,amsthm}
\usepackage[margin=1.1in]{geometry}
\usepackage{booktabs}
\newtheorem{theorem}{Theorem}
\newtheorem{remark}{Remark}
\title{Improved lower bounds for the $3x+1$ problem via the\\ Krasikov--Lagarias inequalities at depths $3^{13}$ through $3^{20}$}
\author{Martien de Jong\thanks{With extensive machine assistance (Claude/Jengo); all results
independently verifiable from the deposited certificates.}}
\date{July 2026 (draft)}
\begin{document}
\maketitle

\begin{abstract}
Let $\pi_1(x)$ denote the number of integers $n\le x$ whose forward orbit under the $3x+1$ map
contains $1$. Krasikov and Lagarias (2003) proved $\pi_1(x)\ge x^{0.84}$ for sufficiently large
$x$ by exhibiting feasible solutions of a linear-program family $L_k^{NT}(\lambda)$ at depth
$k=11$; this has remained the strongest unconditional lower bound (a December 2025 preprint,
arXiv:2512.13760, still cites $x^{0.84}$ as the record). We solve $L_k^{NT}(\lambda)$
at depths $k=13,15,17,19,20,21$ (up to $3{,}486{,}784{,}401$ constraints) and verify explicit
rational certificates in exact integer arithmetic, obtaining
\[
\pi_1(x)\;\ge\; x^{0.9184}\qquad\text{for all sufficiently large }x,
\]
with intermediate verified exponents $0.8624$, $0.8805$, $0.8953$, $0.9069$, $0.9146$; the
$k=21$ exponent was pre-registered as a prediction ($\gamma\approx0.918$, frozen 2026-07-16)
before the computation. No new mathematics is
required beyond \cite{KL}: Theorem~2.2 there converts any feasible solution into the density
bound with exponent $\log_2\lambda$. Our contribution is computational: a monotone (nonlinear
Perron) iteration replaces the interior-point solver, making depths $3^{19}$--$3^{20}$ feasible
on a desktop machine (and $3^{20}$ classes at $k=21$ via memmapped strictly-sequential sweeps
on an external drive), and every certificate is re-verified with all weights replaced by strict
rational lower bounds. We further observe an exact scalar identity governing the family
(the ``min-loss identity''), which predicts the feasibility edges to three decimals and
suggests the method attains $x^{1-\varepsilon}$ for every $\varepsilon>0$.
\end{abstract}

\section{The system}
Fix $k\ge 2$ and let $[3^k]=\{m \bmod 3^k:\ m\equiv 2 \bmod 3\}$, $\alpha=\log_2 3$.
Following \cite[\S2]{KL}, the linear program $L_k^{NT}(\lambda)$ has principal variables
$c^m\ge 1$ for $m\in[3^k]$, auxiliaries
$\bar c^{\,t}=\min(c^{t},c^{t+3^{k-1}},c^{t+2\cdot 3^{k-1}})$ for $t\in[3^{k-1}]$, and
constraints
\begin{align}
m\equiv 2\ (9):\quad & c^m \le c^{4m}\lambda^{-2}+\bar c^{\,(4m-2)/3}\lambda^{\alpha-2},\\
m\equiv 5\ (9):\quad & c^m \le c^{4m}\lambda^{-2},\\
m\equiv 8\ (9):\quad & c^m \le c^{4m}\lambda^{-2}+\bar c^{\,(2m-1)/3}\lambda^{\alpha-1},
\end{align}
with indices taken mod $3^k$ resp.\ $3^{k-1}$. By \cite[Thm.~2.2]{KL}, a feasible positive
solution at parameter $\lambda$ yields $\pi_a(x)\ge \Delta_1 c^m \lambda^{y}$ for every
$a\not\equiv 0 \bmod 3$, hence a density exponent $\gamma=\log_2\lambda$.

\section{Method and results}
The constraint map $F_\lambda$ (right-hand sides, with auxiliaries substituted) is monotone,
positively homogeneous and concave; feasibility is equivalent to its nonlinear Perron root
being $\ge 1$, located by power iteration and bisection in $\lambda$. The implementation
reproduces the published values at $k=9$ ($\gamma=0.8168$; cf.\ $0.81$ in \cite{AL}) and
$k=11$ ($\gamma=0.8418$; cf.\ $0.84$ in \cite{KL}). At greater depths:

\begin{center}
\begin{tabular}{lrrrr}
\toprule
$k$ & classes & $\lambda_0$ (rational) & min.\ slack & verified exponent\\
\midrule
13 & 531{,}441 & $1818/1000$ & $1.000279$ & $0.8624$\\
15 & 4{,}782{,}969 & $1841/1000$ & $1.000307$ & $0.8805$\\
17 & 43{,}046{,}721 & $93/50$ & $1.000504$ & $0.8953$\\
19 & 387{,}420{,}489 & $15/8$ & $1.000939$ & $0.9069$\\
20 & 1{,}162{,}261{,}467 & $1885/1000$ & $1.000194$ & $0.9146$\\
21 & 3{,}486{,}784{,}401 & $1890/1000$ & $1.000187$ & $\mathbf{0.9184}$\\
\bottomrule
\end{tabular}
\end{center}

\begin{theorem}
$\pi_1(x)\ \ge\ x^{0.9184}$ for all sufficiently large $x$.
\end{theorem}

\emph{Verification.} For each depth the deposited integer vector $C$ (scaled Perron vector)
is checked against all constraints with $\lambda_0$ rational and every weight replaced by a
strict rational lower bound (80-digit decimal floors minus one ulp, common denominator
$10^{18}$); comparisons are exact integer arithmetic. Zero violations at all six depths
($5.1\times 10^9$ constraints in total, of which $3.49\times 10^9$ at $k=21$ alone). Certificates and a standalone {\tt numpy}-only
verifier are deposited with this note. Three independent cross-checks strengthen the claim:
(i) a second, independently written implementation (different indexing and $\lambda$
parametrization) reproduces the feasibility edges at overlapping depths to $3$--$4$ decimals
($k{=}13$: $0.8620$; $k{=}15$: $0.8801$; $k{=}17$: $0.8950$);
(ii) the $k=12$ certificate ($\gamma = 213/250 = 0.852$, already above the previous record)
has been machine-verified in \emph{Lean~4}: the coefficient lower bounds reduce to integer
power inequalities (e.g.\ $p_a^{250}\le 2^{250Q-426}$) and all $177{,}147$ class inequalities
are checked by compiled kernel evaluation, with no floating point anywhere;
(iii) both implementations were calibrated against all published anchor values
($k=2,9,11$).

\section{The min-loss identity and the ceiling}
Summing the edge equalities over all classes (the maps $m\mapsto 4m$ and the two branch maps
are class bijections) collapses the system to a scalar law: at the Perron edge,
\[
1\;=\;\lambda^{-2}+\frac{q}{3}\left(\lambda^{\alpha-2}+\lambda^{\alpha-1}\right),
\qquad q \;=\; \frac{3\sum_t \bar c^{\,t}}{\sum_m c^m}.
\]
(Verified to $10^{-8}$ at $k=5,\dots,11$.) Note $q=1$ forces $\lambda=2$, i.e.\ $\gamma=1$:
the entire gap between this method and full density is the ``min-loss'' $1-q$. Measured
$q$ rises $0.888\ (k{=}5)\to 0.9705\ (k{=}19)$, driven by equalization of the refinement
triples of the Perron vector; the within-vector heterogeneity profile decays geometrically
in the trit position with bulk factor $\approx 0.84$, matching the subleading spectral ratio
of the edge-linearized operator. If this contraction admits a uniform bound below $1$
(supported by all measurements; we leave it as a conjecture), then $q\to 1$ and the method
attains $\pi_1(x)\ge x^{1-\varepsilon}$ for every $\varepsilon>0$, as hoped in \cite{KL}.

\begin{remark}
Depth $k=21$ was completed after the first draft of this note: the pre-registered
prediction $\gamma\approx0.918$ (frozen before the computation) was confirmed at
$\gamma=0.9184$, on a desktop by memmapped strictly-sequential sweeps. Depths
$k=25$--$30$ (cloud scale) would discriminate between the extrapolations
$\gamma_\infty\approx0.95$ (ceiling) and $\gamma_\infty=1$ (density); all
flow-fraction measurements to date favor the latter.
\end{remark}

\begin{thebibliography}{9}
\bibitem{KL} I.~Krasikov, J.~C.~Lagarias, \emph{Bounds for the 3x+1 problem using difference
inequalities}, Acta Arith.\ 109 (2003), 237--258; arXiv:math/0205002.
\bibitem{AL} D.~Applegate, J.~C.~Lagarias, \emph{Density bounds for the 3x+1 problem II:
Krasikov inequalities}, Math.\ Comp.\ 64 (1995), 427--438.
\end{thebibliography}
\end{document}
